\section{Database record reference}\label{app:record-reference}

A record is the information for one calculation, written as text on one line.
This appendix is the quick reference; Chapter~\ref{sec:jsonl} explains how to
read a database without programming.

\subsection{Entry order and presence}

\CJKQC{} writes the entries in this fixed order, leaving out information that
is not available:

\begin{verbatim}
status, error, name, source_path, program, program_version, command,
electronic_energy, scf_energy, charge, multiplicity, geometry,
vibrational_frequencies, gradient, dipole_moment, mulliken_charges,
rotational_symmetry_number, components, correction
\end{verbatim}

The entries may appear in any order, and unknown entries are ignored. Missing
information is left out completely; you will not see \texttt{null}. Entry
names use underscores and do not change between versions.

\subsection{Entries}

\begin{longtable}{@{}L{6.0cm}L{2.7cm}L{5.3cm}@{}}
\toprule
\textbf{Entry} & \textbf{Kind of value} & \textbf{Meaning} \\
\midrule
\endhead
\fld{status} & text & \fld{ok}, \fld{failed}, or \fld{truncated}. A missing
entry means \fld{ok}. \\
\fld{error} & text & Reason for a failure; never empty when
\fld{status} is \fld{failed}. \\
\fld{name} & text & File stem, exactly as written. The name used to match
corrections and monomer references. \\
\fld{source_path} & text & Full path of the original output file. \\
\fld{program} & text & Program name: \fld{gaussian}, \fld{orca}, \fld{xtb},
\fld{mrcc}, \fld{xyz}. \\
\fld{program_version} & text & Version as printed by the program. \\
\fld{command} & text & Command line as echoed by the program. \\
\fld{electronic_energy} & number & Final electronic energy, hartree. \\
\fld{scf_energy} & number & Converged SCF energy, hartree. Used only by
extrapolation. \\
\fld{charge} & whole number & Net charge in elementary charges. \\
\fld{multiplicity} & whole number & Spin multiplicity, at least 1. \\
\fld{geometry} & two lists & \fld{symbols} and \fld{coordinates}; coordinates
in \AA{}, listed atom by atom. \\
\fld{vibrational_frequencies} & list & Genuine modes in cm$^{-1}$; a negative
value means imaginary. \\
\fld{gradient} & list & Cartesian gradient, hartree/\AA{}, three values per
atom. \\
\fld{dipole_moment} & list & $[x,y,z]$ in atomic units ($e\,a_0$). \\
\fld{mulliken_charges} & list & Per-atom charges in elementary charges. \\
\fld{rotational_symmetry_number} & whole number & $\sigma$, at least 1. \\
\fld{components} & code and count & Each component code with its count, in
alphabetical order. \\
\fld{correction} & group of values & Attached correction; see below. \\
\bottomrule
\end{longtable}

\subsection{Geometry entries}

\begin{verbatim}
"geometry": {
  "symbols": ["N", "H"],
  "coordinates": [0.0, 0.0, 0.0, 0.0, 0.0, 1.01]
}
\end{verbatim}

\fld{symbols} holds the element names, with the first letter capital. The
\fld{coordinates} list has exactly three numbers per symbol, in the same atom
order: $x,y,z$ for the first atom, then $x,y,z$ for the second, and so on. The
number of coordinate values must be exactly three times the number of symbols.

\subsection{Components entry}

\begin{verbatim}
"components": {"am": 2, "sa": 3}
\end{verbatim}

Component codes are made of letters (A--Z and a--z), and upper and lower case
matter. Counts are at least 1. The codes are written in alphabetical order.

\subsection{Correction entries}

\begin{verbatim}
"correction": {
  "energy": -76.4109,
  "program": "orca",
  "program_version": "6.0.0",
  "commands": ["! wB97x-3c def2-TZVPP Sp"],
  "kind": "sp"
}
\end{verbatim}

\begin{longtable}{@{}L{4.2cm}L{2.3cm}L{7.0cm}@{}}
\toprule
\textbf{Entry} & \textbf{Kind of value} & \textbf{Meaning} \\
\midrule
\endhead
\fld{energy} & number & Correction energy in hartree. \\
\fld{program} & text & Program name of the correction calculation. \\
\fld{program_version} & text & Version of the correction calculation. \\
\fld{commands} & list & Command lines that produced the correction, in order. \\
\fld{kind} & text & How the energy was obtained: \texttt{sp}, \texttt{cbs},
\texttt{cps}, or \texttt{cbs-cps}. \\
\bottomrule
\end{longtable}

\subsection{When a record is rejected}

A record is rejected (and, during collection, replaced by a failed record) if
any of the following holds:

\begin{itemize}
  \item a number is not a finite value (it is infinitely large, or not a
        number at all);
  \item the geometry has no symbols, or its coordinate count is not exactly
        three times the symbol count;
  \item a per-atom list (\fld{gradient}, \fld{mulliken_charges}) disagrees
        with the geometry's atom count when a geometry is present;
  \item a \fld{failed} record has no error message;
  \item \fld{multiplicity} or \fld{rotational_symmetry_number} is below 1;
  \item \fld{components} is empty, a code is empty or contains non-letters, or
        a count is below 1;
  \item a correction has a non-finite energy or an empty \fld{kind}.
\end{itemize}

\subsection{Example records}

A minimal record (the actual file stores this on one line):

\begin{verbatim}
{"status":"ok","electronic_energy":-1.5}
\end{verbatim}

A record with a geometry and an attached correction, shown spread over several
lines for readability:

\begin{verbatim}
{
  "status": "ok",
  "name": "1dma",
  "program": "xtb",
  "electronic_energy": -38.51,
  "geometry": {
    "symbols": ["N", "H"],
    "coordinates": [0.0, 0.0, 0.0, 0.0, 0.0, 1.01]
  },
  "components": {"dma": 1},
  "correction": {
    "energy": -76.4109,
    "program": "orca",
    "program_version": "6.0.0",
    "commands": ["! wB97x-3c def2-TZVPP Sp"],
    "kind": "sp"
  }
}
\end{verbatim}
